\documentclass[noanswers, addpoints, 12pt]{exam}
\usepackage{tikz}
\usetikzlibrary{positioning,shapes.multipart,arrows.meta,shadows}
\usepackage[utf8]{inputenc}
\usepackage[french]{babel}
\usepackage[T1]{fontenc}
\usepackage{xcolor}
\usepackage{listings}
\usepackage{graphicx}
\usepackage[most]{tcolorbox}

\definecolor{codegreen}{rgb}{0,0.6,0}
\definecolor{codegray}{rgb}{0.5,0.5,0.5}
\definecolor{codepurple}{rgb}{0.58,0,0.82}
\definecolor{backcolour}{rgb}{0.95,0.95,0.92}
 
\lstdefinestyle{mystyle}{
    backgroundcolor=\color{backcolour},   
    commentstyle=\color{codegreen},
    keywordstyle=\color{magenta},
    numberstyle=\tiny\color{codegray},
    stringstyle=\color{codepurple},
    basicstyle=\footnotesize,
    breakatwhitespace=false,         
    breaklines=true,                 
    captionpos=b,                    
    keepspaces=true,                 
    numbers=left,                    
    numbersep=5pt,                  
    showspaces=false,                
    showstringspaces=false,
    showtabs=false,                  
    tabsize=2
}
\lstset{style=mystyle, language=Python}

\title{CC1 Term NSI - Pile, File et rappels de 1er}
\author{Babash - NSI Tani Malandi et Johanne Bakalara}
\date{Septembre 2024, Durée: 1h30}


\begin{document}
\runningfooter{}{\footnotesize Babash - NSI Tani Malandi et Johanne Bakalara \textperiodcentered{} CC BY-SA 4.0 \textperiodcentered{} page \thepage}{}
\firstpagefooter{}{\footnotesize Babash - NSI Tani Malandi et Johanne Bakalara \textperiodcentered{} CC BY-SA 4.0 \textperiodcentered{} page \thepage}{}

\maketitle

\pagestyle{headandfoot}
\thispagestyle{headandfoot}
\firstpageheader{Nom : \\  Prénom :\\ Classe :}{}{}


\section*{Consignes:}
\begin{center} 
\fbox{\parbox{\linewidth}{
\begin{center}
Le devoir est à rédiger sur une feuille libre.\\
La clarté de la rédaction et les explications font partie intégrante de la notation. Une réponse non rédigée sans explication sera considérée fausse.\\
Prenez le temps d'encadrer les réponses.
\end{center}


}}
\end{center}

\vsword{Note}
\hsword{Note}

\section*{Barème :}
\begin{center}
  \par\bigskip
  \gradetable[h]
\end{center}
\begin{questions}


\section{Cours : }
\question[4] Questions de cours :
\begin{parts}
    \part Quelle est la différence principale entre un tuple et une liste ?
    \part Soit le code suivant :
        \begin{lstlisting}[language=Python]
repertoire = [{"Prenom": "Alice", "Numero": "0600000001"},
              {"Prenom": "Bob", "Numero": "0600000002"}]
        \end{lstlisting}
        Quelle instruction faut-il écrire pour obtenir le numéro de Bob ?
    \part Quelle structure de données correspond à un ordonnancement LIFO ("Last In First Out") ?
    \part Un employé supermarché est en train de ranger des boites de lait de coco dans un rayon, les unes derrière les autres. L'employé remarque que la première boite mise en rayon sera la dernière prise par les clients. Cette situation correspond-elle à une pile ou à une file ?

\end{parts}

\section{Compréhension de code :}
On considère un tableau d'entiers tab~\textbf{dont les éléments sont des 0 ou
des 1}.
\begin{lstlisting}[language=Python]
def mystere(tab):
    i = 0
    j = len(tab) - 1
    while i != j :
        if tab[i] == 0:
            i = i + 1
        else :
            valeur = tab[j] 
            tab[j] = 1
            tab[i] = valeur
            j= j - 1
    return tab
\end{lstlisting}
\question[2] Donner dans le tableau l'évolution des variables i, j et tab à chaque tour de boucle pour un appel de la fonction mystere([0, 1, 0, 1, 0, 1]).\\

\question[2] Réfléchir sur la fonction :
\begin{parts}
    \part Quelle est la valeur de retour de la fonction mystere pour les appels suivants :
    \begin{itemize}
     \item mystere([0, 0, 0])
     \item mystere([])
     \item mystere([0, 0, 1])
     \item mystere([1, 1, 1, 1])
    \end{itemize}
    
    \part Quel nom pourrait-on donner à la fonction pour mieux décrire son utilité ?
\end{parts}
\section{Rédaction d'un code simple :}
\question[4] Écrire une fonction moyenne qui prend en paramètre un tableau non vide de nombres
flottants et qui renvoie la moyenne des valeurs du tableau. Les tableaux seront représentés sous forme de liste Python.\\
Exemples :
\begin{lstlisting}[language=Python]
>>> moyenne([1.0])
1.0
>>> moyenne([1.0, 2.0, 4.0])
2.3333333333333335
\end{lstlisting}


\section{Exercice Type Bac pile/file}
\textit{Thème abordé : structures de données : les piles}\\\\
On cherche à obtenir un mélange d'une liste comportant un nombre pair d'éléments.
Dans cet exercice, on notera N le nombre d'éléments de la liste à mélanger.
La méthode de mélange utilisée dans cette partie est inspirée d'un mélange de jeux
de cartes :
\begin{itemize}
    \item On sépare la liste en deux piles :
    \begin{itemize}
        \item à gauche, la première pile contient les N/2 premiers éléments de la liste ;
        \item à droite, la deuxième pile contient les N/2 derniers éléments de la liste.
    \end{itemize}
    \item On crée une liste vide.
    \item On prend alors le sommet de la pile de gauche et on le met en début de liste.
    \item On prend ensuite le sommet de la pile de droite que l'on ajoute à la liste et ainsi
de suite jusqu'à ce que les piles soient vides.
\end{itemize}

Par exemple, si on applique cette méthode de mélange à la liste
['V','D','R','3','7','10'], on obtient pour le partage de la liste en 2 piles :
\begin{center}
    \begin{tikzpicture}[font=\small,cell/.style={draw,minimum width=2cm,minimum height=0.6cm,font=\ttfamily}]
\node at (0,2.1) {Pile gauche}; \node[cell] at (0,1.5) {'R'}; \node[cell] at (0,0.9) {'D'}; \node[cell] at (0,0.3) {'V'};
\node at (4,2.1) {Pile droite}; \node[cell] at (4,1.5) {'10'}; \node[cell] at (4,0.9) {'7'}; \node[cell] at (4,0.3) {'3'};
\end{tikzpicture}

\end{center}
La nouvelle liste à la fin du mélange sera donc ['R','10','D','7','V','3'].

\question[1] Que devient la liste ['7','8','9','10','V','D','R','A'] si on lui applique
cette méthode de mélange ?

On considère qu'on dispose de la structure de données de type pile, munie des
seules instructions suivantes :
\begin{itemize}
    \item p = Pile() : crée une pile vide nommée p
    \item p.est\_vide() : renvoie Vrai si la liste est vide, Faux sinon
    \item p.empiler(e) : ajoute l'élément e dans la pile
    \item e = p.depiler() : retire le dernier élément ajouté dans la pile et le retourne (et
l'affecte à la variable e)
\end{itemize}

\question[2] Recopier et compléter le code de la fonction suivante qui transforme une liste
en pile. Le dernier élément de la liste doit se retrouver en haut de la pile.
\begin{lstlisting}[language=Python]
def liste_vers_pile(L):
    '''prend en paramètre une liste et renvoie une pile'''
    N = len(L)
    p_temp = Pile()
    for i in range(N):
        # completer ici
    return ……… # completer le retour de la fonction
\end{lstlisting}

\question[2] On considère la fonction partage suivante qui partage une liste en deux piles quels affichages obtient-on à l'écran lors de l'exécution de l'instruction :
partage([1,2,3,4,5,6]) ?\\\\
Aide : Lors de sa mise au point et pour aider au débuggage, des appels à la fonction affichage\_pile ont été insérés. La fonction affichage\_pile(p) affiche la pile p à l'écran verticalement, c'est la seule instruction d'affichage de la fonction.

\begin{lstlisting}[language=Python]
def partage(L):
    N = len(L)
    p_gauche = Pile()
    p_droite = Pile()
    for i in range(N // 2):
        p_gauche.empiler(L[i])
    for i in range(N // 2, N):
        p_droite.empiler(L[i])
    affichage_pile(p_gauche)
    affichage_pile(p_droite)
    return p_gauche, p_droite
\end{lstlisting}

\question[1] Dans un cas général et en vous appuyant sur une séquence de schémas, expliquer en quelques lignes comment fusionner deux piles p\_gauche et p\_droite pour former une liste L en alternant un à un les éléments de la pile p\_gauche et de la pile p\_droite.

\question[2] Écrire une fonction fusion(p1, p2) qui renvoie une liste construite à partir
des deux piles p1 et p2.

\end{questions}

\end{document}

